% ============================================================================
% Fichier  : partie1/P01_C01_audit_controle_auditologie.tex
% Titre    : Auditer, Contrôler, Donner une Assurance : du Métier à l'Auditologie
% Auteur   : MINKA MI NGUIDJOÏ Thierry Emmanuel
% Version  : 1.0 -- octobre 2026
% ============================================================================
% Sources : CRIPP 2017 de l'IIA (dossier Normes, définition de l'audit
% interne, paraphrasée) ; modèle des trois lignes (IIA, 2020) ; guide d'audit
% des SI du CHAI (2014), typologie des audits ; INTOSAI GUID 5100 ; article
% A0 « Vers l'auditologie » (HAL hal-05701007) pour la définition de
% l'auditologie et les axiomes ; loi n° 2010/012 (texte anglais officiel,
% vérifié pour le cours GCI). Dates de l'histoire de l'ISACA (1969, 1978,
% 1996) : présentations publiques de l'ISACA.
% ============================================================================

\chapter[Auditer, Contrôler, Donner une Assurance]{Auditer, Contrôler, Donner une Assurance : du Métier à l'Auditologie}
\label{chap:audit-auditologie}
\chaptermark{Du métier à l'auditologie}

\epigraph{%
    « L'auditeur ne dit pas que tout va bien.\\
    Il dit ce qu'il sait, comment il le sait, et ce qu'il ignore. »%
}{Adage de la profession, reformulé}

\niveauChapitre{Fondamental}{%
    Maîtrise du Chapitre~\ref{chap:prologue} (impression, constatation,
    opinion~; réponse \NEK{})%
}

\begin{boxobjectifs}
\begin{itemize}
    \item Distinguer audit, contrôle, inspection, conseil et certification.
    \item Définir l'assurance raisonnable, l'assurance limitée et leurs
          limites.
    \item Situer l'audit interne, l'audit externe et l'audit public, et les
          rôles du modèle des trois lignes.
    \item Retracer les grandes étapes de l'audit informatique, de l'audit
          « autour de l'ordinateur » à l'audit continu.
    \item Présenter les grands types d'audit des systèmes d'information.
    \item Présenter l'audit comme réduction de l'incertitude sur un système
          sociotechnique, et ce que cette lecture change pour la pratique.
\end{itemize}
\end{boxobjectifs}

\bigskip

% ============================================================================
\section*{Question directrice : qu'achète-t-on quand on paie un audit~?}
% ============================================================================

Après l'échec de la première mission (Chapitre~\ref{chap:prologue}),
M.~ETEME, directeur général de LiZbiZ, pose une question simple au chef de
mission d'AUDIX CONSEIL~: «~Si je vous paie, pouvez-vous me garantir qu'il
n'y a plus de fraude dans mes entrepôts~?~» M.~BEYALA répond non. Il peut
garantir une démarche, des preuves et une opinion motivée sur ce qu'il aura
examiné. Il ne peut pas garantir l'absence de toute fraude.

Cette réponse déçoit souvent les dirigeants. Elle est pourtant au cœur du
métier. Ce chapitre définit ce qu'est un audit, ce qu'il garantit et ce
qu'il ne garantit pas (sections~\ref{sec:c01-notions}
et~\ref{sec:c01-assurance}), qui audite (section~\ref{sec:c01-acteurs}),
d'où vient l'audit informatique (section~\ref{sec:c01-histoire}), quels sont
ses grands types (section~\ref{sec:c01-types}), et pourquoi ce cours le lit
comme une \emph{réduction de l'incertitude} (section~\ref{sec:c01-auditologie}).

\separateur

% ============================================================================
\section{Audit, Contrôle, Inspection, Conseil}
\label{sec:c01-notions}
% ============================================================================

\subsection{Cinq activités voisines}

Dans une organisation, plusieurs activités examinent le fonctionnement des
systèmes. Elles se ressemblent, mais elles ne posent pas la même question et
ne produisent pas le même livrable.

\begin{boxdef}
\termecle{Audit}\index{audit}~: examen méthodique, indépendant et documenté,
qui confronte une situation à des critères connus à l'avance, et conduit à
une opinion motivée sur cette situation, fondée sur des preuves.

\medskip
\termecle{Contrôle}\index{contrôle}~: dispositif ou action, intégré au
fonctionnement de l'organisation, qui vise à maîtriser une activité au fil
de l'eau (vérifier une saisie, approuver un paiement, rapprocher deux
comptes). Le contrôle est exercé par l'organisation elle-même.

\medskip
\termecle{Inspection}\index{inspection}~: vérification ponctuelle de la
conformité d'un objet ou d'une opération à une règle, souvent par une
autorité de tutelle, et souvent à visée de sanction.

\medskip
\termecle{Conseil}\index{conseil}~: accompagnement qui aide l'organisation à
concevoir ou à améliorer un dispositif. Le conseiller propose~; il ne donne
pas d'opinion indépendante sur ce qu'il a contribué à concevoir.

\medskip
\termecle{Certification}\index{certification}~: attestation, par un
organisme tiers accrédité, qu'un système de management est conforme aux
exigences d'une norme, sur un périmètre et pour une durée donnés.
\end{boxdef}

\begin{tableaucours}{L{2.4cm}L{3.5cm}L{3.0cm}L{3.3cm}}{%
    \tableauHeader{Activité} &
    \tableauHeader{Question posée} &
    \tableauHeader{Qui l'exerce} &
    \tableauHeader{Livrable}
}
Audit & La situation est-elle conforme aux critères, et dans quelle mesure
le sait-on~? & Auditeur indépendant de l'activité & Rapport et opinion
motivée \\
Contrôle & L'opération en cours est-elle correcte~? & Les opérationnels,
la hiérarchie & Une opération validée ou rejetée \\
Inspection & La règle est-elle respectée ici et maintenant~? & Une autorité
de tutelle & Procès-verbal, sanction éventuelle \\
Conseil & Comment faire mieux~? & Consultant, expert & Recommandations,
dispositif conçu \\
Certification & Le système satisfait-il la norme~? & Organisme accrédité &
Certificat \\
\end{tableaucours}
\vspace{-0.8em}
\begin{center}{\small\itshape Tableau~1.1. Cinq activités d'examen à ne pas
confondre}\end{center}

\subsection{L'audit interne selon l'IIA}

L'Institut des auditeurs internes (\sigle{IIA}) définit l'audit interne
comme une activité indépendante et objective, qui donne à une organisation
une assurance sur le degré de maîtrise de ses opérations, lui apporte des
conseils pour les améliorer et contribue à créer de la valeur. Il y parvient
en évaluant, par une approche systématique et méthodique, ses processus de
gouvernance, de management des risques et de contrôle~\cite{iia2017cripp}.

Trois mots de cette définition méritent l'attention. \emph{Assurance}~:
l'audit produit un degré de confiance, pas une garantie. \emph{Conseil}~:
l'audit interne peut aussi conseiller, à condition de ne pas auditer ensuite
ce qu'il a conçu (Chapitre~\ref{chap:normes}). \emph{Systématique}~:
l'audit suit une méthode, que la Partie~\ref{part:methode} enseigne.

\begin{boxalert}
\textbf{Violation typique.} Le contrôleur interne de LiZbiZ rédige la
procédure de désactivation des comptes, la met en place, puis est chargé
«~d'auditer~» son application.

\smallskip
\textbf{Règle~:} on n'audite pas ce qu'on a conçu ou ce qu'on exécute. Le
contrôle relève de l'organisation~; l'audit doit en être indépendant.
\end{boxalert}

% ============================================================================
\section{Ce que l'Audit Garantit, et ce qu'il ne Garantit pas}
\label{sec:c01-assurance}
% ============================================================================

\subsection{Assurance raisonnable et assurance limitée}

\begin{boxdef}
\termecle{Assurance raisonnable}\index{assurance raisonnable}~: niveau élevé,
mais non absolu, de confiance que l'auditeur exprime sous une forme
positive («~nous estimons que le dispositif est efficace~»), après des
travaux étendus.

\medskip
\termecle{Assurance limitée}\index{assurance limitée}~: niveau de confiance
plus faible, exprimé sous une forme négative («~nous n'avons rien relevé
qui laisse penser que\ldots~»), après des travaux moins étendus, comme une
revue documentaire.
\end{boxdef}

L'assurance n'est jamais absolue, pour quatre raisons qui tiennent à la
nature même de l'audit.

\begin{enumerate}
    \item \textbf{L'auditeur échantillonne.} Il n'examine pas toutes les
          opérations~: une anomalie peut se trouver hors de l'échantillon.
    \item \textbf{Les contrôles peuvent être contournés}, notamment par
          collusion ou par un dirigeant qui passe outre.
    \item \textbf{L'auditeur juge.} Le choix des critères, des tests et de
          la matérialité repose sur un jugement professionnel.
    \item \textbf{L'auditeur ne voit pas tout.} Une partie de l'information
          est absente, cachée ou contradictoire. C'est là que la réponse
          \NEK{} prend sa place (section~\ref{sec:c01-auditologie}).
\end{enumerate}

\begin{boiteRemarque}{Pourquoi ne pas tout vérifier~?}
Parce que le coût serait démesuré, et que le résultat ne serait pas
meilleur~: vérifier les 2,4~millions de mouvements de stock annuels de
LiZbiZ ne dirait rien des mouvements qui n'ont pas été enregistrés.
L'auditeur cherche donc les bons tests et le bon échantillon
(Chapitre~\ref{chap:preuve}), pas l'exhaustivité.
\end{boiteRemarque}

\subsection{La réponse à M. ETEME}

\begin{boxscenario}{Opération LIZBIZ. La lettre de mission}
\textbf{Faits.} Le directeur général demande une garantie d'absence de
fraude.

\medskip
\textbf{Réponse du chef de mission.} La lettre de mission proposée par
AUDIX CONSEIL précise~:
\begin{itemize}
    \item l'objectif~: donner une assurance raisonnable sur l'efficacité des
          contrôles d'accès et des contrôles de gestion des stocks de
          l'application intégrée, sur la période de janvier à avril~2026~;
    \item le périmètre~: le centre de données de Douala et quatre agences
          sur douze, choisies sur la base des écarts d'inventaire~;
    \item ce que la mission ne fera pas~: elle ne recherche pas
          systématiquement la fraude, et ne remplace pas une enquête si des
          indices sérieux apparaissent.
\end{itemize}

\textbf{Analyse} (confiance élevée). La lettre transforme une attente
irréaliste (« plus de fraude ») en un objectif vérifiable. Elle protège
l'auditeur, et elle protège aussi le directeur général, qui sait ce qu'il
achète.
\end{boxscenario}

% ============================================================================
\section{Qui Audite~? Interne, Externe, Public}
\label{sec:c01-acteurs}
% ============================================================================

\subsection{Trois familles d'auditeurs}

\begin{tableaucours}{L{2.6cm}L{3.8cm}L{2.9cm}L{2.9cm}}{%
    \tableauHeader{Famille} &
    \tableauHeader{Rôle} &
    \tableauHeader{Commanditaire} &
    \tableauHeader{Exemple}
}
Audit interne & Assurance et conseil sur la maîtrise des opérations &
Organe de gouvernance (comité d'audit) & Direction de l'audit interne de la
BANQUE SANAGA \\
Audit externe légal & Certification des comptes & Assemblée des associés,
la loi & Le commissaire aux comptes de LiZbiZ \\
Audit externe contractuel & Mission ponctuelle sur un sujet défini & La
direction & AUDIX CONSEIL chez LiZbiZ \\
Audit public & Contrôle de l'emploi des fonds publics & La loi, la
juridiction & La Chambre des comptes à l'ANPP \\
\end{tableaucours}
\vspace{-0.8em}
\begin{center}{\small\itshape Tableau~1.2. Les familles d'auditeurs}\end{center}

Dans le secteur public, l'audit externe relève des institutions supérieures
de contrôle des finances publiques, qui appliquent les normes de
l'INTOSAI. Le GUID~5100 de l'INTOSAI fournit le cadre général de la
vérification des systèmes d'information~\cite{intosai2019guid5100}. Au
Cameroun, la Chambre des comptes de la Cour suprême exerce le contrôle
juridictionnel des comptes publics~; le Contrôle Supérieur de l'État exerce
le contrôle administratif des services publics.

\subsection{Le modèle des trois lignes}

L'IIA a publié en 2020 le \termecle{modèle des trois lignes}\index{modèle des
trois lignes}~\cite{iia2020troislignes}. Il répartit les rôles entre
l'organe de gouvernance, le management et l'audit interne.

\begin{tableaucours}{L{2.8cm}L{5.0cm}L{4.4cm}}{%
    \tableauHeader{Acteur} &
    \tableauHeader{Rôle} &
    \tableauHeader{Chez LiZbiZ}
}
Organe de gouvernance & Fixe la direction, supervise, rend compte aux
parties prenantes & Conseil d'administration, comité d'audit \\
Première ligne & Porte et maîtrise les risques dans ses activités &
Responsable informatique, chefs d'agence, logistique \\
Deuxième ligne & Expertise et surveillance des risques & Contrôleuse
interne, responsable de la sécurité \\
Troisième ligne & Assurance indépendante & Audit interne (à créer) ou
cabinet mandaté \\
Prestataires externes & Assurance complémentaire & Commissaire aux comptes,
AUDIX CONSEIL \\
\end{tableaucours}
\vspace{-0.8em}
\begin{center}{\small\itshape Tableau~1.3. Le modèle des trois lignes
appliqué à LiZbiZ}\end{center}

\begin{boxmetier}{Contrôleur interne}
\textbf{Se situer en deuxième ligne}
\begin{enumerate}
    \item Concevoir et documenter les contrôles avec la première ligne,
          sans les exécuter à sa place.
    \item Surveiller leur fonctionnement par des contrôles de second niveau
          (revue d'échantillons, indicateurs).
    \item Remonter les anomalies à la direction et au comité d'audit, sans
          filtre.
    \item Laisser l'audit, interne ou externe, évaluer le dispositif de façon
          indépendante~: le contrôleur interne ne s'audite pas lui-même.
\end{enumerate}
\end{boxmetier}

% ============================================================================
\section{D'où Vient l'Audit Informatique}
\label{sec:c01-histoire}
% ============================================================================

\subsection{Trois manières d'auditer avec l'ordinateur}

Lorsque l'informatique entre dans les services comptables, dans les années
1960, les auditeurs l'abordent d'abord comme une boîte noire. Trois
approches se succèdent, puis coexistent.

\begin{tableaucours}{L{3.2cm}L{5.0cm}L{4.0cm}}{%
    \tableauHeader{Approche} &
    \tableauHeader{Principe} &
    \tableauHeader{Limite}
}
Autour de l'ordinateur & On rapproche les entrées et les sorties sans
examiner les traitements & Aveugle aux traitements automatisés et aux
volumes \\
À travers l'ordinateur & On teste les contrôles programmés et la logique
des traitements & Demande des compétences techniques \\
Avec l'ordinateur & On utilise des outils pour analyser l'ensemble des
données (techniques d'audit assistées par ordinateur) & Demande des
extractions fiables \\
\end{tableaucours}
\vspace{-0.8em}
\begin{center}{\small\itshape Tableau~1.4. Les approches de l'audit
informatique}\end{center}

\subsection{Une profession qui s'organise}

D'après les présentations publiques de l'ISACA, une association
d'auditeurs du traitement électronique des données est fondée en 1969~; elle
deviendra l'ISACA. La certification d'auditeur des systèmes d'information
(\sigle{CISA}) apparaît en 1978. Le référentiel COBIT paraît en 1996. En
2002, la loi Sarbanes-Oxley~\cite{sox2002} impose aux sociétés cotées aux
États-Unis de rendre compte de l'efficacité de leur contrôle interne sur
l'information financière~: les contrôles informatiques deviennent un objet
d'audit à part entière. Le Chapitre~\ref{chap:normes} présentera les normes
qui encadrent aujourd'hui la profession.

\subsection{Vers l'audit continu}

Avec les volumes de données actuels, l'audit ponctuel montre ses limites~:
une anomalie peut durer des mois entre deux missions. L'\termecle{audit
continu}\index{audit continu} analyse les données et les contrôles de façon
répétée, voire en permanence. Le Chapitre~\ref{chap:audit-continu} y est
consacré.

% ============================================================================
\section{Les Grands Types d'Audit des SI}
\label{sec:c01-types}
% ============================================================================

Le guide d'audit des systèmes d'information du Comité d'harmonisation de
l'audit interne (\sigle{CHAI}) distingue plusieurs thèmes
d'audit~\cite{chai2014guide}. Le cours en retient sept, chacun traité dans
la Partie~\ref{part:pratique}.

\begin{tableaucours}{L{4.0cm}L{5.4cm}c}{%
    \tableauHeader{Type d'audit} &
    \tableauHeader{Question centrale} &
    \tableauHeader{Chapitre}
}
Pilotage et fonction informatique & La fonction informatique sert-elle les
objectifs de l'organisation~? & \ref{chap:audit-gouvernance} \\
Application en service & L'application produit-elle des données exactes et
complètes~? & \ref{chap:audit-applications} \\
Projets et fonction « études » & Le projet livrera-t-il ce qui est
attendu, dans les délais et le budget~? & \ref{chap:audit-projets} \\
Sécurité & L'information est-elle protégée en confidentialité, intégrité,
disponibilité~? & \ref{chap:audit-securite} \\
Production, support et parc & L'exploitation est-elle fiable et
maîtrisée~? & \ref{chap:audit-continuite} \\
Continuité et tiers & Le service survivra-t-il à un sinistre ou à la
défaillance d'un prestataire~? & \ref{chap:audit-continuite} \\
Missions généralistes & Quelle place prend l'informatique dans un audit
financier ou de conformité~? & \ref{chap:audit-applications} \\
\end{tableaucours}
\vspace{-0.8em}
\begin{center}{\small\itshape Tableau~1.5. Les types d'audit des SI (d'après
le guide du CHAI)}\end{center}

\begin{boxscenario}{Opération LIZBIZ. Quel type d'audit~?}
\textbf{Faits.} Écart de stock de 186~millions FCFA, réserve du
commissaire aux comptes, réponses incohérentes sur le contrôle d'accès.

\medskip
\textbf{Analyse} (confiance élevée). La demande initiale («~tout le
système d'information~») n'est pas un type d'audit. La seconde mission
combine deux types~: un audit d'\emph{application en service}, centré sur
la gestion des stocks de l'application intégrée, et un audit de
\emph{sécurité}, limité au contrôle d'accès (domaine~D). Ce découpage
détermine les référentiels, les compétences et la durée.
\end{boxscenario}

% ============================================================================
\section{Le Cadre Camerounais en un Coup d'Œil}
\label{sec:c01-cadre}
% ============================================================================

Au Cameroun, plusieurs textes rendent l'audit ou le contrôle interne
obligatoires pour certaines organisations. Ils seront étudiés au
Chapitre~\ref{chap:droit}~; en voici les principaux.

\begin{itemize}
    \item La loi n°~2010/012 relative à la cybersécurité et à la
          cybercriminalité prévoit un audit de sécurité obligatoire et
          périodique des systèmes d'information de certaines organisations
          (articles~13 et~32)~\cite{loi2010012}.
    \item Le règlement COBAC R-2016/04 fixe les exigences minimales du
          contrôle interne des établissements de crédit de la
          CEMAC~\cite{cobac2016r04}.
    \item Le droit OHADA des sociétés commerciales organise la mission du
          commissaire aux comptes et, pour les sociétés cotées, le comité
          d'audit~\cite{ohada2014auscgie}.
\end{itemize}

\begin{boxjuris}
\textbf{\loicam{n°~2010/012}, article~26, alinéa~1} (version anglaise
officielle)~:

\medskip
\textit{« Operators of information systems shall take every technical and
administrative measure to ensure the security of services offered. To
this end, they shall have standardized systems enabling them to at all
times identify, assess, process or manage any risk relating to the
security of the information systems of the services provided directly or
indirectly. »}

\medskip
\textbf{Portée pour l'audit.} L'exploitant d'un système d'information doit
disposer de dispositifs normalisés pour identifier, évaluer et traiter en
permanence les risques de sécurité. Un auditeur peut donc vérifier
l'existence et le fonctionnement de ces dispositifs~: c'est un critère
d'audit d'origine légale, qui s'impose à LiZbiZ comme à la BANQUE SANAGA.
\end{boxjuris}

% ============================================================================
\section{L'Audit comme Réduction de l'Incertitude}
\label{sec:c01-auditologie}
% ============================================================================

\subsection{Un constat~: un métier sans socle scientifique}

L'audit des SI s'apprend par compagnonnage, se codifie dans des référentiels
professionnels et se pratique par jugement d'expert. Ces pratiques sont
précieuses, mais, comme le soutient l'article fondateur de ce
cours~\cite{minka2026auditologie}, il manque à la discipline un socle
scientifique unificateur~: des axiomes explicites, des propositions
réfutables et des mesures communes pour son objet central.

\subsection{Une proposition~: l'auditologie}

\begin{boxdef}
\termecle{Mission auditologique}\index{mission auditologique}~: ensemble
formé par cinq éléments~: le \emph{système audité}~$\mathcal{S}$ (personnes,
technologies, flux d'information, règles)~; l'\emph{auditeur}~$\mathcal{A}$,
ou l'équipe d'audit, avec ses limites de traitement, de mémoire et
d'attention~; le \emph{canal d'information}~$\mathcal{I}$ par lequel les
preuves parviennent à l'auditeur~; la \emph{structure stratégique}~$\mathcal{G}$
des relations entre auditeur, audité et tiers~; la \emph{trajectoire}~$\mathcal{T}$
du système au fil des missions.

\medskip
\termecle{Auditologie}\index{auditologie}~: science appliquée de la
réduction de l'incertitude sur les missions
auditologiques~\cite{minka2026auditologie}.
\end{boxdef}

L'article pose six axiomes ontologiques, présentés en détail au
Chapitre~\ref{chap:auditologie}. Trois d'entre eux éclairent dès maintenant
la pratique.

\begin{itemize}
    \item[] \textbf{Axiome~1 (Observabilité partielle).} L'auditeur n'a accès
          qu'à une partie des états pertinents du système audité, et cette
          partie est elle-même partiellement opaque.
    \item[] \textbf{Axiome~4 (Interaction stratégique).} L'auditeur et l'audité
          poursuivent des intérêts distincts~; leurs actions se comprennent
          comme des coups dans un jeu à information imparfaite.
    \item[] \textbf{Axiome~6 (Opposabilité).} Toute conclusion d'une mission
          doit pouvoir être défendue devant un tiers qui a autorité pour la
          contester.
\end{itemize}

{\small\itshape Axiomes reformulés d'après~\cite{minka2026auditologie}, avec
la numérotation de l'article~; les six axiomes figurent au
Chapitre~\ref{chap:auditologie}.}

\subsection{Ce que cette lecture change pour la pratique}

\begin{tableaucours}{L{4.4cm}L{7.8cm}}{%
    \tableauHeader{Lecture classique} &
    \tableauHeader{Lecture auditologique}
}
L'audit vérifie la conformité & L'audit réduit l'incertitude~; la
conformité à un référentiel est un \emph{choix d'objectif}, pas un
fondement \\
Une question appelle OUI ou NON & L'observabilité partielle impose une
troisième réponse, \NEK{} (Chapitre~\ref{chap:preuve}) \\
L'indépendance est une règle déontologique & L'indépendance se déduit de
l'interaction stratégique et de l'opposabilité \\
Le taux de conformité résume un domaine & Un domaine se lit par trois
mesures, $\mu$, $U$ et $H$ (Chapitre~\ref{chap:conclure}) \\
Une grille longue est une grille complète & La grille doit respecter les
limites cognitives de l'auditeur~: on décompose
(Chapitre~\ref{chap:decomposition}) \\
\end{tableaucours}
\vspace{-0.8em}
\begin{center}{\small\itshape Tableau~1.6. Deux lectures de
l'audit}\end{center}

\begin{boiteRemarque}{L'auditologie remplace-t-elle les normes professionnelles~?}
Non. Elle ne prescrit pas le «~bon audit~» et n'invalide pas les cadres
existants~\cite{minka2026auditologie}. Les normes de l'IIA, de l'ISACA ou de
l'INTOSAI disent ce que l'auditeur doit faire~; l'auditologie cherche à
expliquer \emph{pourquoi} les méthodes fonctionnent ou échouent, et à le
vérifier par l'expérience. Ce cours enseigne les deux.
\end{boiteRemarque}

\begin{boxmetier}{Auditeur SI}
\textbf{Appliquer la lecture auditologique dès la première réunion}
\begin{enumerate}
    \item Annoncer à l'audité que la réponse \NEK{} est admise, et qu'elle
          sera documentée comme telle.
    \item Distinguer, dans le plan de travail, ce qui sera prouvé, ce qui
          sera seulement déclaré, et ce qui restera inconnu.
    \item Prévoir, pour chaque domaine, le calcul de $\mu$, $U$ et $H$.
    \item Rédiger chaque conclusion de façon qu'un tiers puisse la refaire à
          partir des papiers de travail.
\end{enumerate}
\end{boxmetier}

\separateur

\begin{boxresume}
\textbf{Ce qu'il faut retenir de ce chapitre~:}
\begin{itemize}
    \item \textbf{Audit, contrôle, inspection, conseil, certification} sont
          cinq activités distinctes~; l'audit est indépendant et conclut sur
          des preuves au regard de critères.
    \item \textbf{L'assurance n'est jamais absolue}~: l'auditeur
          échantillonne, juge, ne voit pas tout, et les contrôles peuvent
          être contournés.
    \item \textbf{Trois familles d'auditeurs}~: interne, externe (légal ou
          contractuel), public~; et \textbf{trois lignes} de maîtrise des
          risques sous l'organe de gouvernance.
    \item \textbf{L'audit informatique} est passé de l'audit autour de
          l'ordinateur à l'audit avec l'ordinateur, puis à l'audit continu.
    \item \textbf{Un type d'audit} se choisit~: pilotage, application,
          projets, sécurité, production, continuité et tiers.
    \item \textbf{L'auditologie} lit l'audit comme une réduction de
          l'incertitude sur un système partiellement observable, peuplé
          d'acteurs stratégiques~; elle justifie \NEK{} et les trois mesures
          $\mu$, $U$, $H$.
\end{itemize}
\end{boxresume}

\bigskip

\begin{boxexo}[title={\ding{45}\quad Exercice~1.1. Quelle activité~? \niveaubase}]
Pour chacune des situations suivantes, indiquez s'il s'agit d'un audit,
d'un contrôle, d'une inspection, d'un conseil ou d'une certification, et
justifiez.
\begin{enumerate}[(a)]
    \item Le chef d'agence de Garoua vérifie chaque soir que les bons de
          sortie de l'entrepôt sont signés.
    \item La COBAC se rend dans une agence de la BANQUE SANAGA pour vérifier
          le respect d'une règle de conservation des pièces.
    \item AUDIX CONSEIL donne une opinion sur l'efficacité des contrôles
          d'accès de LiZbiZ.
    \item Un cabinet aide l'ANPP à rédiger sa politique de sécurité.
\end{enumerate}
\end{boxexo}

\begin{boxexo}[title={\ding{45}\quad Exercice~1.2. Les trois lignes de la BANQUE SANAGA \niveaumoyen}]
La BANQUE SANAGA compte une DSI, un RSSI rattaché au directeur des
risques, une direction de l'audit interne dirigée par Mme~EKAMBI, et un
commissaire aux comptes.
\begin{enumerate}[(a)]
    \item Placez chacun de ces acteurs dans le modèle des trois lignes.
    \item Mme~EKAMBI a participé, en 2025, à la rédaction de la procédure de
          gestion des changements. Quel problème cela pose-t-il pour un
          audit de cette procédure~?
    \item Proposez une solution conforme à l'indépendance de l'audit.
\end{enumerate}
\end{boxexo}

\begin{boxexo}[title={\ding{45}\quad Exercice~1.3. Une lettre de mission \niveauexpert}]
La Chambre des comptes annonce un audit des systèmes d'information de
l'ANPP.
\begin{enumerate}[(a)]
    \item Rédigez l'objectif, le périmètre et les limites de la mission, sur
          le modèle de la lettre d'AUDIX CONSEIL.
    \item Indiquez le niveau d'assurance visé (raisonnable ou limitée) et
          justifiez.
    \item Identifiez deux situations dans lesquelles l'équipe devra répondre
          \NEK{}, et expliquez ce que ces réponses apprendront à la Chambre.
    \item Montrez en quoi l'axiome d'opposabilité s'applique avec une force
          particulière à un audit public.
\end{enumerate}
\end{boxexo}

\bigskip

\begin{boxinfo}
\textbf{Transition.} L'audit évalue un dispositif~: le contrôle interne.
Avant d'apprendre à l'auditer, il faut savoir ce qu'il est, de quoi il se
compose et comment on le conçoit. C'est l'objet du
Chapitre~\ref{chap:controle-interne}.
\end{boxinfo}
